\documentclass[11pt]{book}
\usepackage[margin=0.8in]{geometry}
\usepackage{amsmath,amssymb}
\usepackage{array,longtable,enumitem}
\usepackage{hyperref}
\setlist{nosep}
\title{Algebra I Full-Year Computational Course}
\author{Student Book and Teacher Answer Key}
\date{}
\begin{document}
\maketitle
\tableofcontents
\clearpage
\chapter{Algebra I Full-Year Computational Course}

\section{Student Book and Teacher Answer Key}

\textbf{Course design:} full-year Algebra I, unit by unit, with computational practice only. This book intentionally avoids word problems and focuses on symbolic, numeric, tabular, and graph-ready algebra skills.

\subsection{How to Use This Book}

Each unit contains:

\begin{enumerate}
\item \textbf{Learning targets} aligned to Algebra I topics.
\item \textbf{Formula sheet} for quick reference.
\item \textbf{Sample problems} with worked solutions.
\item \textbf{Classwork practice} for guided or independent class use.
\item \textbf{Homework practice} for additional repetition.
\item \textbf{Teacher answer key} for every classwork and homework item.
\end{enumerate}

\bigskip\hrule\bigskip

\chapter{Year-at-a-Glance Pacing}

\begin{longtable}{|p{0.12\textwidth}|p{0.58\textwidth}|p{0.18\textwidth}|}
\hline
Unit & Topic & Suggested Time \\ \hline
1 & Real Numbers, Properties, and Precision & 2 weeks \\ \hline
2 & Exponents, Radicals, and Rational Exponents & 3 weeks \\ \hline
3 & Expressions, Structure, and Equivalent Forms & 3 weeks \\ \hline
4 & Linear Equations, Inequalities, and Absolute Value & 4 weeks \\ \hline
5 & Coordinate Graphing and Linear Functions & 4 weeks \\ \hline
6 & Systems of Equations and Inequalities & 4 weeks \\ \hline
7 & Functions, Notation, Domain, Range, and Transformations & 4 weeks \\ \hline
8 & Polynomials and Factoring & 4 weeks \\ \hline
9 & Quadratic Functions and Equations & 5 weeks \\ \hline
10 & Exponential Functions and Sequences & 3 weeks \\ \hline
11 & Rational Expressions, Radicals, and Inverses & 3 weeks \\ \hline
12 & Statistics, Data Displays, and Linear Models & 3 weeks \\ \hline
13 & Cumulative Review and Final Exam Practice & 2 weeks \\ \hline
\end{longtable}

\bigskip\hrule\bigskip

\chapter{Coursewide Formula Sheets}

\section{Number and Quantity}

\begin{itemize}
\item Natural numbers: \(1,2,3,\ldots\)
\item Whole numbers: \(0,1,2,3,\ldots\)
\item Integers: \(\ldots,-2,-1,0,1,2,\ldots\)
\item Rational numbers: numbers that can be written as \(\frac{a}{b}\), where \(a,b\) are integers and \(b\ne0\)
\item Irrational numbers: non-repeating, non-terminating decimals that cannot be written as \(\frac{a}{b}\)
\item Absolute value: \(|a|\) is the distance from \(a\) to 0
\item Percent error: \(\frac{|\text{approximate}-\text{exact}|}{|\text{exact}|}\cdot100\%\)
\end{itemize}

\section{Exponents and Radicals}

\begin{itemize}
\item \(a^m\cdot a^n=a^{m+n}\)
\item \(\frac{a^m}{a^n}=a^{m-n}\), \(a\ne0\)
\item \((a^m)^n=a^{mn}\)
\item \((ab)^n=a^nb^n\)
\item \(\left(\frac{a}{b}\right)^n=\frac{a^n}{b^n}\), \(b\ne0\)
\item \(a^0=1\), \(a\ne0\)
\item \(a^{-n}=\frac{1}{a^n}\), \(a\ne0\)
\item \(a^{1/n}=\sqrt[n]{a}\)
\item \(a^{m/n}=\sqrt[n]{a^m}=(\sqrt[n]{a})^m\)
\item \(\sqrt{ab}=\sqrt a\sqrt b\)
\item \(\sqrt{\frac{a}{b}}=\frac{\sqrt a}{\sqrt b}\), \(b>0\)
\end{itemize}

\section{Linear Equations and Inequalities}

\begin{itemize}
\item Slope: \(m=\frac{y_2-y_1}{x_2-x_1}\)
\item Slope-intercept form: \(y=mx+b\)
\item Point-slope form: \(y-y_1=m(x-x_1)\)
\item Standard form: \(Ax+By=C\)
\item Horizontal line: \(y=c\)
\item Vertical line: \(x=c\)
\item Inequality rule: multiplying or dividing by a negative number reverses the inequality sign.
\item Absolute value equations:
\begin{itemize}
\item \(|x|=a\Rightarrow x=a\) or \(x=-a\), if \(a\ge0\)
\item \(|x|<a\Rightarrow -a<x<a\)
\item \(|x|>a\Rightarrow x<-a\) or \(x>a\)
\end{itemize}
\end{itemize}

\section{Systems}

\begin{itemize}
\item Substitution: solve one equation for a variable and substitute.
\item Elimination: add or subtract equations to eliminate a variable.
\item One solution: lines intersect once.
\item No solution: parallel lines.
\item Infinitely many solutions: same line.
\end{itemize}

\section{Functions}

\begin{itemize}
\item Function notation: \(f(x)\) means the output when input is \(x\).
\item Domain: set of allowed inputs.
\item Range: set of possible outputs.
\item Average rate of change: \(\frac{f(b)-f(a)}{b-a}\)
\item Transformation forms:
\begin{itemize}
\item \(f(x)+k\): vertical shift
\item \(f(x+k)\): horizontal shift left \(k\)
\item \(kf(x)\): vertical stretch/compression/reflection
\item \(f(kx)\): horizontal compression/stretch/reflection
\end{itemize}
\end{itemize}

\section{Polynomials and Factoring}

\begin{itemize}
\item Standard form: terms in descending powers.
\item Degree: greatest exponent of the variable.
\item Distributive property: \(a(b+c)=ab+ac\)
\item Difference of squares: \(a^2-b^2=(a-b)(a+b)\)
\item Perfect square trinomials:
\begin{itemize}
\item \(a^2+2ab+b^2=(a+b)^2\)
\item \(a^2-2ab+b^2=(a-b)^2\)
\end{itemize}
\item Factoring \(x^2+bx+c\): find numbers with sum \(b\) and product \(c\).
\item Factoring \(ax^2+bx+c\): use grouping or another valid method.
\end{itemize}

\section{Quadratics}

\begin{itemize}
\item Standard form: \(f(x)=ax^2+bx+c\)
\item Vertex form: \(f(x)=a(x-h)^2+k\)
\item Factored form: \(f(x)=a(x-r_1)(x-r_2)\)
\item Axis of symmetry: \(x=-\frac{b}{2a}\)
\item Vertex: \(\left(-\frac{b}{2a}, f\left(-\frac{b}{2a}\right)\right)\)
\item Quadratic formula: \(x=\frac{-b\pm\sqrt{b^2-4ac}}{2a}\)
\item Discriminant: \(D=b^2-4ac\)
\begin{itemize}
\item \(D>0\): two real solutions
\item \(D=0\): one real solution
\item \(D<0\): two complex solutions
\end{itemize}
\item Completing the square: \(x^2+bx=(x+\frac b2)^2-(\frac b2)^2\)
\end{itemize}

\section{Exponential Functions and Sequences}

\begin{itemize}
\item Exponential form: \(f(x)=ab^x\)
\item Growth: \(b>1\)
\item Decay: \(0<b<1\)
\item Arithmetic sequence: \(a_n=a_1+(n-1)d\)
\item Geometric sequence: \(a_n=a_1r^{n-1}\)
\item Recursive arithmetic: \(a_n=a_{n-1}+d\)
\item Recursive geometric: \(a_n=ra_{n-1}\)
\end{itemize}

\section{Statistics}

\begin{itemize}
\item Mean: \(\bar{x}=\frac{\text{sum of values}}{\text{number of values}}\)
\item Median: middle value after ordering.
\item Range: maximum minus minimum.
\item Interquartile range: \(IQR=Q_3-Q_1\)
\item Mean absolute deviation: average distance from the mean.
\item Residual: observed \(y\)-value minus predicted \(y\)-value.
\item Correlation coefficient \(r\): measures strength and direction of linear association.
\end{itemize}

\bigskip\hrule\bigskip

\chapter{Unit 1: Real Numbers, Properties, and Precision}

\section{Learning Targets}

\begin{itemize}
\item Classify numbers as rational, irrational, integer, whole, or natural.
\item Use properties of rational and irrational numbers.
\item Apply order of operations and algebraic properties.
\item Round values and choose appropriate precision.
\end{itemize}

\section{Unit Formula Sheet}

\begin{itemize}
\item Rational + rational = rational
\item Rational \(\cdot\) rational = rational
\item Rational + irrational = irrational
\item Nonzero rational \(\cdot\) irrational = irrational
\item Commutative: \(a+b=b+a\), \(ab=ba\)
\item Associative: \((a+b)+c=a+(b+c)\), \((ab)c=a(bc)\)
\item Distributive: \(a(b+c)=ab+ac\)
\end{itemize}

\section{Sample Problems}

\subsection{Sample 1}

Classify \(-\frac{8}{4}\).

\[-\frac{8}{4}=-2\]

It is rational and an integer.

\subsection{Sample 2}

Simplify \(3[8-(2+1)^2]+5\).

\[(2+1)^2=9\]
\[8-9=-1\]
\[3(-1)+5=2\]

Answer: \(2\)

\subsection{Sample 3}

State whether \(5+\sqrt7\) is rational or irrational.

A rational number plus an irrational number is irrational.

Answer: irrational

\section{Classwork}

\begin{enumerate}
\item Classify \(\sqrt{25}\) as rational or irrational.
\item Classify \(\sqrt{18}\) as rational or irrational.
\item Simplify \(4(3+2)-7\).
\item Simplify \(2^3+6\div3\).
\item Simplify \(-3(4-9)+2\).
\item Simplify \(\frac{3}{4}+\frac{5}{6}\).
\item Simplify \(\frac{7}{8}-\frac{1}{3}\).
\item Simplify \(-5+|-12|\).
\item Determine whether \(\sqrt2\cdot 6\) is rational or irrational.
\item Round \(18.7462\) to the nearest hundredth.
\end{enumerate}

\section{Homework}

\begin{enumerate}
\item Classify \(-9\) as natural, whole, integer, rational, or irrational.
\item Classify \(0\) as whole, integer, rational, or irrational.
\item Simplify \(6^2-4(5)\).
\item Simplify \(18\div3+2^4\).
\item Simplify \(-2[7-(3^2)]\).
\item Simplify \(\frac{2}{5}+\frac{7}{10}\).
\item Simplify \(\frac{11}{12}-\frac{1}{8}\).
\item Simplify \(|-15|-|-6|\).
\item Determine whether \(\sqrt5+8\) is rational or irrational.
\item Round \(3.14159\) to the nearest thousandth.
\end{enumerate}

\section{Teacher Answer Key}

\subsection{Classwork Answers}

\begin{enumerate}
\item rational
\item irrational
\item \(13\)
\item \(10\)
\item \(17\)
\item \(\frac{19}{12}\)
\item \(\frac{13}{24}\)
\item \(7\)
\item irrational
\item \(18.75\)
\end{enumerate}

\subsection{Homework Answers}

\begin{enumerate}
\item integer and rational
\item whole, integer, rational
\item \(16\)
\item \(22\)
\item \(4\)
\item \(\frac{11}{10}\)
\item \(\frac{19}{24}\)
\item \(9\)
\item irrational
\item \(3.142\)
\end{enumerate}

\bigskip\hrule\bigskip

\chapter{Unit 2: Exponents, Radicals, and Rational Exponents}

\section{Learning Targets}

\begin{itemize}
\item Apply exponent properties.
\item Rewrite radicals using rational exponents.
\item Rewrite rational exponents using radicals.
\item Simplify radicals and expressions with rational exponents.
\end{itemize}

\section{Unit Formula Sheet}

\begin{itemize}
\item \(a^{m/n}=\sqrt[n]{a^m}\)
\item \(\sqrt[n]{a}=a^{1/n}\)
\item \(a^{-n}=\frac1{a^n}\)
\item \(a^0=1\), \(a\ne0\)
\item \(\sqrt{a^2}=|a|\)
\end{itemize}

\section{Sample Problems}

\subsection{Sample 1}

Simplify \(x^4\cdot x^7\).

\[x^4\cdot x^7=x^{11}\]

\subsection{Sample 2}

Rewrite \(\sqrt[3]{y^5}\) using rational exponents.

\[\sqrt[3]{y^5}=y^{5/3}\]

\subsection{Sample 3}

Simplify \(16^{3/4}\).

\[16^{3/4}=(\sqrt[4]{16})^3=2^3=8\]

\section{Classwork}

\begin{enumerate}
\item Simplify \(x^3x^8\).
\item Simplify \(\frac{a^{10}}{a^4}\).
\item Simplify \((m^2)^5\).
\item Simplify \((3x^2)(4x^5)\).
\item Simplify \(y^{-3}\).
\item Simplify \(9^{1/2}\).
\item Simplify \(27^{1/3}\).
\item Simplify \(81^{3/4}\).
\item Rewrite \(\sqrt{x^7}\) using rational exponents.
\item Rewrite \(p^{5/2}\) using radicals.
\end{enumerate}

\section{Homework}

\begin{enumerate}
\item Simplify \(b^6b^9\).
\item Simplify \(\frac{r^{12}}{r^5}\).
\item Simplify \((n^4)^3\).
\item Simplify \((5q^3)(2q^4)\).
\item Simplify \(z^{-5}\).
\item Simplify \(25^{1/2}\).
\item Simplify \(64^{1/3}\).
\item Simplify \(32^{2/5}\).
\item Rewrite \(\sqrt[4]{a^9}\) using rational exponents.
\item Rewrite \(t^{7/3}\) using radicals.
\end{enumerate}

\section{Teacher Answer Key}

\subsection{Classwork Answers}

\begin{enumerate}
\item \(x^{11}\)
\item \(a^6\)
\item \(m^{10}\)
\item \(12x^7\)
\item \(\frac1{y^3}\)
\item \(3\)
\item \(3\)
\item \(27\)
\item \(x^{7/2}\)
\item \(\sqrt{p^5}\) or \(p^2\sqrt p\)
\end{enumerate}

\subsection{Homework Answers}

\begin{enumerate}
\item \(b^{15}\)
\item \(r^7\)
\item \(n^{12}\)
\item \(10q^7\)
\item \(\frac1{z^5}\)
\item \(5\)
\item \(4\)
\item \(4\)
\item \(a^{9/4}\)
\item \(\sqrt[3]{t^7}\)
\end{enumerate}

\bigskip\hrule\bigskip

\chapter{Unit 3: Expressions, Structure, and Equivalent Forms}

\section{Learning Targets}

\begin{itemize}
\item Identify terms, coefficients, constants, and factors.
\item Combine like terms.
\item Use distribution and factoring to create equivalent expressions.
\item Rewrite expressions to reveal structure.
\end{itemize}

\section{Unit Formula Sheet}

\begin{itemize}
\item Like terms have identical variable parts.
\item Coefficient: numerical factor multiplying a variable term.
\item Constant: term with no variable.
\item Factor common factor: \(ab+ac=a(b+c)\)
\item Standard linear expression: \(ax+b\)
\end{itemize}

\section{Sample Problems}

\subsection{Sample 1}

Simplify \(4x+7-2x+9\).

\[4x-2x+7+9=2x+16\]

\subsection{Sample 2}

Factor \(12x+18\).

\[12x+18=6(2x+3)\]

\subsection{Sample 3}

Rewrite \(3(x+4)+2(x-5)\) in standard form.

\[3x+12+2x-10=5x+2\]

\section{Classwork}

\begin{enumerate}
\item Simplify \(7x+3x\).
\item Simplify \(9a-4a+6\).
\item Simplify \(5(2x+3)\).
\item Simplify \(-2(4y-7)\).
\item Simplify \(3(x-2)+4(x+5)\).
\item Simplify \(8m+2-5m+11\).
\item Factor \(10x+15\).
\item Factor \(14a-21\).
\item Factor \(6x^2+9x\).
\item Identify the coefficient of \(-8p^3\).
\end{enumerate}

\section{Homework}

\begin{enumerate}
\item Simplify \(12y-5y\).
\item Simplify \(8b+9-3b\).
\item Simplify \(7(3x-1)\).
\item Simplify \(-5(2n+6)\).
\item Simplify \(6(a+1)-2(a-9)\).
\item Simplify \(11q-4+2q+10\).
\item Factor \(18x+24\).
\item Factor \(20m-35\).
\item Factor \(12r^2-8r\).
\item Identify the constant in \(9x^2-4x+11\).
\end{enumerate}

\section{Teacher Answer Key}

\subsection{Classwork Answers}

\begin{enumerate}
\item \(10x\)
\item \(5a+6\)
\item \(10x+15\)
\item \(-8y+14\)
\item \(7x+14\)
\item \(3m+13\)
\item \(5(2x+3)\)
\item \(7(2a-3)\)
\item \(3x(2x+3)\)
\item \(-8\)
\end{enumerate}

\subsection{Homework Answers}

\begin{enumerate}
\item \(7y\)
\item \(5b+9\)
\item \(21x-7\)
\item \(-10n-30\)
\item \(4a+24\)
\item \(13q+6\)
\item \(6(3x+4)\)
\item \(5(4m-7)\)
\item \(4r(3r-2)\)
\item \(11\)
\end{enumerate}

\bigskip\hrule\bigskip

\chapter{Unit 4: Linear Equations, Inequalities, and Absolute Value}

\section{Learning Targets}

\begin{itemize}
\item Solve one-step, two-step, and multi-step linear equations.
\item Solve equations with variables on both sides.
\item Solve and graph linear inequalities.
\item Solve absolute value equations and inequalities.
\item Rearrange literal equations.
\end{itemize}

\section{Unit Formula Sheet}

\begin{itemize}
\item Equation solving goal: isolate the variable.
\item Use inverse operations.
\item Keep equations balanced by doing the same operation to both sides.
\item Inequality reversal: if multiplying or dividing by a negative, reverse \(<,>,\le,\ge\).
\item \(|ax+b|=c\Rightarrow ax+b=c\) or \(ax+b=-c\), if \(c\ge0\).
\end{itemize}

\section{Sample Problems}

\subsection{Sample 1}

Solve \(3x+5=20\).

\[3x=15\]
\[x=5\]

\subsection{Sample 2}

Solve \(-2x+7\le15\).

\[-2x\le8\]
\[x\ge-4\]

\subsection{Sample 3}

Solve \(|x-3|=8\).

\[x-3=8 \quad \text{or} \quad x-3=-8\]
\[x=11 \quad \text{or} \quad x=-5\]

\section{Classwork}

\begin{enumerate}
\item Solve \(x+8=15\).
\item Solve \(4x=28\).
\item Solve \(3x-7=11\).
\item Solve \(5(x+2)=35\).
\item Solve \(7x-4=3x+12\).
\item Solve \(\frac{x}{3}+5=9\).
\item Solve \(2x-1<9\).
\item Solve \(-3x+6\ge15\).
\item Solve \(|x+4|=10\).
\item Solve \(A=lw\) for \(w\).
\end{enumerate}

\section{Homework}

\begin{enumerate}
\item Solve \(x-9=14\).
\item Solve \(-6x=42\).
\item Solve \(2x+13=31\).
\item Solve \(-4(x-3)=20\).
\item Solve \(9x+1=4x+26\).
\item Solve \(\frac{x}{5}-2=6\).
\item Solve \(5x+3\le18\).
\item Solve \(-2x-7>5\).
\item Solve \(|2x-1|=9\).
\item Solve \(P=2l+2w\) for \(l\).
\end{enumerate}

\section{Teacher Answer Key}

\subsection{Classwork Answers}

\begin{enumerate}
\item \(x=7\)
\item \(x=7\)
\item \(x=6\)
\item \(x=5\)
\item \(x=4\)
\item \(x=12\)
\item \(x<5\)
\item \(x\le-3\)
\item \(x=6\) or \(x=-14\)
\item \(w=\frac{A}{l}\)
\end{enumerate}

\subsection{Homework Answers}

\begin{enumerate}
\item \(x=23\)
\item \(x=-7\)
\item \(x=9\)
\item \(x=-2\)
\item \(x=5\)
\item \(x=40\)
\item \(x\le3\)
\item \(x<-6\)
\item \(x=5\) or \(x=-4\)
\item \(l=\frac{P-2w}{2}\)
\end{enumerate}

\bigskip\hrule\bigskip

\chapter{Unit 5: Coordinate Graphing and Linear Functions}

\section{Learning Targets}

\begin{itemize}
\item Find slope from points, tables, and equations.
\item Graph lines in slope-intercept, point-slope, and standard form.
\item Write equations of lines.
\item Identify parallel and perpendicular slopes.
\end{itemize}

\section{Unit Formula Sheet}

\begin{itemize}
\item \(m=\frac{y_2-y_1}{x_2-x_1}\)
\item \(y=mx+b\)
\item \(y-y_1=m(x-x_1)\)
\item Parallel lines have equal slopes.
\item Perpendicular lines have slopes whose product is \(-1\).
\item \(x\)-intercept: set \(y=0\).
\item \(y\)-intercept: set \(x=0\).
\end{itemize}

\section{Sample Problems}

\subsection{Sample 1}

Find the slope through \((2,5)\) and \((6,13)\).

\[m=\frac{13-5}{6-2}=\frac84=2\]

\subsection{Sample 2}

Write the equation with slope \(3\) and \(y\)-intercept \(-4\).

\[y=3x-4\]

\subsection{Sample 3}

Find the \(x\)-intercept of \(2x+3y=12\).

Set \(y=0\):
\[2x=12\]
\[x=6\]

Answer: \((6,0)\)

\section{Classwork}

\begin{enumerate}
\item Find the slope through \((1,2)\) and \((4,8)\).
\item Find the slope through \((-2,5)\) and \((3,-5)\).
\item Write the equation with slope \(4\) and \(y\)-intercept \(1\).
\item Write the equation with slope \(-2\) and \(y\)-intercept \(7\).
\item Find the slope of \(y=\frac32x-9\).
\item Find the \(y\)-intercept of \(y=-5x+4\).
\item Find the \(x\)-intercept of \(3x+2y=18\).
\item Find the \(y\)-intercept of \(4x-5y=20\).
\item Write a line parallel to \(y=2x-6\) through \((0,5)\).
\item Write a line perpendicular to \(y=\frac12x+3\) through \((0,-1)\).
\end{enumerate}

\section{Homework}

\begin{enumerate}
\item Find the slope through \((0,4)\) and \((2,10)\).
\item Find the slope through \((-3,-1)\) and \((1,7)\).
\item Write the equation with slope \(\frac12\) and \(y\)-intercept \(-3\).
\item Write the equation with slope \(-4\) and \(y\)-intercept \(6\).
\item Find the slope of \(y=-\frac45x+2\).
\item Find the \(y\)-intercept of \(y=7x-8\).
\item Find the \(x\)-intercept of \(5x+y=30\).
\item Find the \(y\)-intercept of \(2x+3y=15\).
\item Write a line parallel to \(y=-3x+1\) through \((0,4)\).
\item Write a line perpendicular to \(y=-\frac14x+2\) through \((0,9)\).
\end{enumerate}

\section{Teacher Answer Key}

\subsection{Classwork Answers}

\begin{enumerate}
\item \(2\)
\item \(-2\)
\item \(y=4x+1\)
\item \(y=-2x+7\)
\item \(\frac32\)
\item \((0,4)\)
\item \((6,0)\)
\item \((0,-4)\)
\item \(y=2x+5\)
\item \(y=-2x-1\)
\end{enumerate}

\subsection{Homework Answers}

\begin{enumerate}
\item \(3\)
\item \(2\)
\item \(y=\frac12x-3\)
\item \(y=-4x+6\)
\item \(-\frac45\)
\item \((0,-8)\)
\item \((6,0)\)
\item \((0,5)\)
\item \(y=-3x+4\)
\item \(y=4x+9\)
\end{enumerate}

\bigskip\hrule\bigskip

\chapter{Unit 6: Systems of Equations and Inequalities}

\section{Learning Targets}

\begin{itemize}
\item Solve linear systems by graphing, substitution, and elimination.
\item Classify systems as one solution, no solution, or infinitely many solutions.
\item Solve systems involving a line and a quadratic.
\item Solve and graph systems of linear inequalities.
\end{itemize}

\section{Unit Formula Sheet}

\begin{itemize}
\item Solution of a system: ordered pair satisfying every equation.
\item Elimination works because adding a multiple of one equation to another preserves the solution set.
\item Linear-quadratic systems: substitute the linear expression into the quadratic equation.
\end{itemize}

\section{Sample Problems}

\subsection{Sample 1}

Solve:

\[y=x+2\]
\[y=3x-4\]

Set equal:
\[x+2=3x-4\]
\[6=2x\]
\[x=3\]
\[y=5\]

Answer: \((3,5)\)

\subsection{Sample 2}

Solve:

\[2x+y=7\]
\[3x-y=8\]

Add:
\[5x=15\]
\[x=3\]
\[2(3)+y=7\Rightarrow y=1\]

Answer: \((3,1)\)

\section{Classwork}

\begin{enumerate}
\item Solve: \(y=x+1\), \(y=2x-3\).
\item Solve: \(y=-x+5\), \(y=x+1\).
\item Solve: \(x+y=9\), \(x-y=1\).
\item Solve: \(2x+y=10\), \(x-y=2\).
\item Solve: \(3x+2y=16\), \(x+2y=8\).
\item Classify: \(y=2x+1\), \(y=2x-4\).
\item Classify: \(2x+y=5\), \(4x+2y=10\).
\item Solve: \(y=x^2\), \(y=4\).
\item Solve: \(y=x^2+1\), \(y=5\).
\item State one solution to the inequality system \(y>x\), \(y<5\).
\end{enumerate}

\section{Homework}

\begin{enumerate}
\item Solve: \(y=x+4\), \(y=3x\).
\item Solve: \(y=-2x+6\), \(y=x-3\).
\item Solve: \(x+y=12\), \(x-y=4\).
\item Solve: \(4x+y=18\), \(x-y=3\).
\item Solve: \(5x+y=19\), \(2x-y=2\).
\item Classify: \(y=-x+2\), \(y=-x+7\).
\item Classify: \(3x-y=6\), \(6x-2y=12\).
\item Solve: \(y=x^2\), \(y=9\).
\item Solve: \(y=x^2-2\), \(y=7\).
\item State one solution to the inequality system \(y<2x\), \(y>0\).
\end{enumerate}

\section{Teacher Answer Key}

\subsection{Classwork Answers}

\begin{enumerate}
\item \((4,5)\)
\item \((2,3)\)
\item \((5,4)\)
\item \((4,2)\)
\item \((4,2)\)
\item no solution
\item infinitely many solutions
\item \((-2,4)\), \((2,4)\)
\item \((-2,5)\), \((2,5)\)
\item Sample: \((1,2)\)
\end{enumerate}

\subsection{Homework Answers}

\begin{enumerate}
\item \((2,6)\)
\item \((3,0)\)
\item \((8,4)\)
\item \((\frac{21}{5},\frac65)\)
\item \((3,4)\)
\item no solution
\item infinitely many solutions
\item \((-3,9)\), \((3,9)\)
\item \((-3,7)\), \((3,7)\)
\item Sample: \((1,1)\)
\end{enumerate}

\bigskip\hrule\bigskip

\chapter{Unit 7: Functions, Notation, Domain, Range, and Transformations}

\section{Learning Targets}

\begin{itemize}
\item Determine whether relations are functions.
\item Evaluate function notation.
\item Find domain and range from sets, tables, and equations.
\item Calculate average rate of change.
\item Identify transformations of functions.
\item Find inverse functions for simple linear functions.
\end{itemize}

\section{Unit Formula Sheet}

\begin{itemize}
\item Function rule: each input has exactly one output.
\item \(f(a)\): substitute \(a\) for \(x\).
\item Average rate of change: \(\frac{f(b)-f(a)}{b-a}\)
\item Inverse of a linear function: replace \(f(x)\) with \(y\), switch \(x\) and \(y\), solve for \(y\).
\end{itemize}

\section{Sample Problems}

\subsection{Sample 1}

If \(f(x)=2x^2-3\), find \(f(4)\).

\[f(4)=2(4)^2-3=32-3=29\]

\subsection{Sample 2}

Find the inverse of \(f(x)=3x-5\).

\[y=3x-5\]
\[x=3y-5\]
\[x+5=3y\]
\[y=\frac{x+5}{3}\]

Answer: \(f^{-1}(x)=\frac{x+5}{3}\)

\section{Classwork}

\begin{enumerate}
\item If \(f(x)=3x+2\), find \(f(5)\).
\item If \(g(x)=x^2-4\), find \(g(-3)\).
\item If \(h(x)=\frac{x}{2}+7\), find \(h(10)\).
\item Find the domain of \(\{(-2,5),(0,1),(4,9)\}\).
\item Find the range of \(\{(-2,5),(0,1),(4,9)\}\).
\item Is \(\{(1,2),(1,3),(4,5)\}\) a function?
\item Find the average rate of change of \(f(x)=x^2\) from \(x=1\) to \(x=4\).
\item Describe the transformation \(f(x)+6\).
\item Describe the transformation \(f(x-3)\).
\item Find the inverse of \(f(x)=2x+8\).
\end{enumerate}

\section{Homework}

\begin{enumerate}
\item If \(f(x)=4x-1\), find \(f(6)\).
\item If \(g(x)=x^2+2x\), find \(g(3)\).
\item If \(h(x)=\frac{x-4}{3}\), find \(h(13)\).
\item Find the domain of \(\{(-1,7),(2,8),(5,8)\}\).
\item Find the range of \(\{(-1,7),(2,8),(5,8)\}\).
\item Is \(\{(2,4),(3,4),(4,4)\}\) a function?
\item Find the average rate of change of \(f(x)=2x+1\) from \(x=-2\) to \(x=3\).
\item Describe the transformation \(f(x)-4\).
\item Describe the transformation \(2f(x)\).
\item Find the inverse of \(f(x)=5x-10\).
\end{enumerate}

\section{Teacher Answer Key}

\subsection{Classwork Answers}

\begin{enumerate}
\item \(17\)
\item \(5\)
\item \(12\)
\item \(\{-2,0,4\}\)
\item \(\{1,5,9\}\)
\item no
\item \(5\)
\item vertical shift up 6
\item horizontal shift right 3
\item \(f^{-1}(x)=\frac{x-8}{2}\)
\end{enumerate}

\subsection{Homework Answers}

\begin{enumerate}
\item \(23\)
\item \(15\)
\item \(3\)
\item \(\{-1,2,5\}\)
\item \(\{7,8\}\)
\item yes
\item \(2\)
\item vertical shift down 4
\item vertical stretch by factor 2
\item \(f^{-1}(x)=\frac{x+10}{5}\)
\end{enumerate}

\bigskip\hrule\bigskip

\chapter{Unit 8: Polynomials and Factoring}

\section{Learning Targets}

\begin{itemize}
\item Add, subtract, and multiply polynomials.
\item Classify polynomials by degree and number of terms.
\item Factor common factors, trinomials, and special products.
\item Use factoring to reveal zeros.
\end{itemize}

\section{Unit Formula Sheet}

\begin{itemize}
\item Monomial: one term.
\item Binomial: two terms.
\item Trinomial: three terms.
\item Degree of polynomial: greatest exponent after simplification.
\item FOIL: \((a+b)(c+d)=ac+ad+bc+bd\)
\item Zero product property: if \(ab=0\), then \(a=0\) or \(b=0\).
\end{itemize}

\section{Sample Problems}

\subsection{Sample 1}

Add \((3x^2+2x-5)+(x^2-7x+9)\).

\[4x^2-5x+4\]

\subsection{Sample 2}

Multiply \((x+4)(x-6)\).

\[x^2-6x+4x-24=x^2-2x-24\]

\subsection{Sample 3}

Factor \(x^2+7x+12\).

Numbers with product \(12\) and sum \(7\): \(3\) and \(4\).

\[x^2+7x+12=(x+3)(x+4)\]

\section{Classwork}

\begin{enumerate}
\item Add \((2x^2+3x+1)+(5x^2-x+4)\).
\item Subtract \((7x^2+2x-8)-(3x^2+5x-1)\).
\item Multiply \(3x(2x^2-4x+5)\).
\item Multiply \((x+2)(x+5)\).
\item Multiply \((x-3)(x+7)\).
\item Multiply \((2x+1)(x-4)\).
\item Factor \(x^2+9x+20\).
\item Factor \(x^2-5x+6\).
\item Factor \(x^2-16\).
\item Factor \(2x^2+7x+3\).
\end{enumerate}

\section{Homework}

\begin{enumerate}
\item Add \((4x^2-2x+6)+(x^2+8x-9)\).
\item Subtract \((9x^2-x+2)-(4x^2+3x+7)\).
\item Multiply \(-2x(5x^2+x-3)\).
\item Multiply \((x+6)(x+1)\).
\item Multiply \((x-8)(x+2)\).
\item Multiply \((3x-2)(x+5)\).
\item Factor \(x^2+8x+15\).
\item Factor \(x^2-7x+10\).
\item Factor \(x^2-25\).
\item Factor \(3x^2+10x+3\).
\end{enumerate}

\section{Teacher Answer Key}

\subsection{Classwork Answers}

\begin{enumerate}
\item \(7x^2+2x+5\)
\item \(4x^2-3x-7\)
\item \(6x^3-12x^2+15x\)
\item \(x^2+7x+10\)
\item \(x^2+4x-21\)
\item \(2x^2-7x-4\)
\item \((x+4)(x+5)\)
\item \((x-2)(x-3)\)
\item \((x-4)(x+4)\)
\item \((2x+1)(x+3)\)
\end{enumerate}

\subsection{Homework Answers}

\begin{enumerate}
\item \(5x^2+6x-3\)
\item \(5x^2-4x-5\)
\item \(-10x^3-2x^2+6x\)
\item \(x^2+7x+6\)
\item \(x^2-6x-16\)
\item \(3x^2+13x-10\)
\item \((x+3)(x+5)\)
\item \((x-5)(x-2)\)
\item \((x-5)(x+5)\)
\item \((3x+1)(x+3)\)
\end{enumerate}

\bigskip\hrule\bigskip

\chapter{Unit 9: Quadratic Functions and Equations}

\section{Learning Targets}

\begin{itemize}
\item Graph quadratic functions and identify key features.
\item Solve quadratic equations by factoring, square roots, completing the square, and quadratic formula.
\item Use the discriminant to classify solutions.
\item Rewrite quadratics in equivalent forms.
\end{itemize}

\section{Unit Formula Sheet}

\begin{itemize}
\item \(f(x)=ax^2+bx+c\)
\item Axis: \(x=-\frac{b}{2a}\)
\item Vertex: substitute the axis value into the function.
\item Factored zeros: \(a(x-r_1)(x-r_2)=0\Rightarrow x=r_1,r_2\)
\item Quadratic formula: \(x=\frac{-b\pm\sqrt{b^2-4ac}}{2a}\)
\item Complete the square: add \((\frac b2)^2\) to both sides after isolating \(x^2+bx\).
\end{itemize}

\section{Sample Problems}

\subsection{Sample 1}

Solve \(x^2-9=0\).

\[x^2=9\]
\[x=\pm3\]

\subsection{Sample 2}

Solve \(x^2+5x+6=0\).

\[(x+2)(x+3)=0\]
\[x=-2,-3\]

\subsection{Sample 3}

Find the vertex of \(f(x)=x^2-4x+1\).

\[x=-\frac{-4}{2(1)}=2\]
\[f(2)=4-8+1=-3\]

Vertex: \((2,-3)\)

\section{Classwork}

\begin{enumerate}
\item Solve \(x^2=49\).
\item Solve \(x^2-25=0\).
\item Solve \(x^2+6x+8=0\).
\item Solve \(x^2-3x-10=0\).
\item Solve \(2x^2+5x+2=0\).
\item Find the axis of \(y=x^2-8x+3\).
\item Find the vertex of \(y=x^2+2x-5\).
\item Find the discriminant of \(x^2-4x+9=0\).
\item Solve \(x^2+4x+1=0\) using the quadratic formula.
\item Complete the square: rewrite \(x^2+10x\) as a squared binomial minus a constant.
\end{enumerate}

\section{Homework}

\begin{enumerate}
\item Solve \(x^2=81\).
\item Solve \(x^2-36=0\).
\item Solve \(x^2+7x+12=0\).
\item Solve \(x^2-x-20=0\).
\item Solve \(3x^2+11x+6=0\).
\item Find the axis of \(y=x^2+10x-2\).
\item Find the vertex of \(y=x^2-6x+4\).
\item Find the discriminant of \(2x^2+3x-5=0\).
\item Solve \(x^2-2x-3=0\) using the quadratic formula.
\item Complete the square: rewrite \(x^2-12x\) as a squared binomial minus a constant.
\end{enumerate}

\section{Teacher Answer Key}

\subsection{Classwork Answers}

\begin{enumerate}
\item \(x=\pm7\)
\item \(x=\pm5\)
\item \(x=-2,-4\)
\item \(x=5,-2\)
\item \(x=-\frac12,-2\)
\item \(x=4\)
\item \((-1,-6)\)
\item \(-20\)
\item \(x=-2\pm\sqrt3\)
\item \((x+5)^2-25\)
\end{enumerate}

\subsection{Homework Answers}

\begin{enumerate}
\item \(x=\pm9\)
\item \(x=\pm6\)
\item \(x=-3,-4\)
\item \(x=5,-4\)
\item \(x=-\frac23,-3\)
\item \(x=-5\)
\item \((3,-5)\)
\item \(49\)
\item \(x=3,-1\)
\item \((x-6)^2-36\)
\end{enumerate}

\bigskip\hrule\bigskip

\chapter{Unit 10: Exponential Functions and Sequences}

\section{Learning Targets}

\begin{itemize}
\item Identify exponential growth and decay.
\item Evaluate exponential functions.
\item Rewrite exponential expressions to show rates.
\item Write arithmetic and geometric sequences explicitly and recursively.
\item Compare linear, quadratic, and exponential patterns from tables.
\end{itemize}

\section{Unit Formula Sheet}

\begin{itemize}
\item Exponential function: \(f(x)=ab^x\)
\item Growth if \(b>1\); decay if \(0<b<1\)
\item Arithmetic explicit: \(a_n=a_1+(n-1)d\)
\item Geometric explicit: \(a_n=a_1r^{n-1}\)
\item Arithmetic recursive: \(a_1=\text{first term}\), \(a_n=a_{n-1}+d\)
\item Geometric recursive: \(a_1=\text{first term}\), \(a_n=ra_{n-1}\)
\end{itemize}

\section{Sample Problems}

\subsection{Sample 1}

Evaluate \(f(x)=3(2)^x\) when \(x=4\).

\[f(4)=3(2)^4=48\]

\subsection{Sample 2}

Find an explicit formula for \(5,8,11,14,\ldots\).

\[a_1=5,\quad d=3\]
\[a_n=5+3(n-1)\]

\subsection{Sample 3}

Find an explicit formula for \(2,10,50,250,\ldots\).

\[a_1=2,\quad r=5\]
\[a_n=2(5)^{n-1}\]

\section{Classwork}

\begin{enumerate}
\item Evaluate \(2^5\).
\item Evaluate \(4(3)^2\).
\item Determine growth or decay: \(y=5(1.2)^x\).
\item Determine growth or decay: \(y=7(0.8)^x\).
\item Find the next three terms: \(4,9,14,19,\ldots\).
\item Find the next three terms: \(3,6,12,24,\ldots\).
\item Write an explicit formula for \(7,10,13,16,\ldots\).
\item Write an explicit formula for \(5,15,45,135,\ldots\).
\item Write a recursive rule for \(2,7,12,17,\ldots\).
\item Write a recursive rule for \(4,20,100,500,\ldots\).
\end{enumerate}

\section{Homework}

\begin{enumerate}
\item Evaluate \(3^4\).
\item Evaluate \(2(5)^3\).
\item Determine growth or decay: \(y=9(1.05)^x\).
\item Determine growth or decay: \(y=12(0.6)^x\).
\item Find the next three terms: \(1,5,9,13,\ldots\).
\item Find the next three terms: \(2,8,32,128,\ldots\).
\item Write an explicit formula for \(6,11,16,21,\ldots\).
\item Write an explicit formula for \(3,12,48,192,\ldots\).
\item Write a recursive rule for \(9,12,15,18,\ldots\).
\item Write a recursive rule for \(6,30,150,750,\ldots\).
\end{enumerate}

\section{Teacher Answer Key}

\subsection{Classwork Answers}

\begin{enumerate}
\item \(32\)
\item \(36\)
\item growth
\item decay
\item \(24,29,34\)
\item \(48,96,192\)
\item \(a_n=7+3(n-1)\)
\item \(a_n=5(3)^{n-1}\)
\item \(a_1=2,\ a_n=a_{n-1}+5\)
\item \(a_1=4,\ a_n=5a_{n-1}\)
\end{enumerate}

\subsection{Homework Answers}

\begin{enumerate}
\item \(81\)
\item \(250\)
\item growth
\item decay
\item \(17,21,25\)
\item \(512,2048,8192\)
\item \(a_n=6+5(n-1)\)
\item \(a_n=3(4)^{n-1}\)
\item \(a_1=9,\ a_n=a_{n-1}+3\)
\item \(a_1=6,\ a_n=5a_{n-1}\)
\end{enumerate}

\bigskip\hrule\bigskip

\chapter{Unit 11: Rational Expressions, Radicals, and Inverses}

\section{Learning Targets}

\begin{itemize}
\item Simplify rational expressions.
\item Multiply and divide rational expressions.
\item Add and subtract rational expressions with common denominators.
\item Simplify radical expressions.
\item Find inverses of linear and simple quadratic-root functions.
\end{itemize}

\section{Unit Formula Sheet}

\begin{itemize}
\item Rational expression restrictions: denominator \(\ne0\).
\item \(\frac{ab}{ac}=\frac bc\), \(a\ne0\), \(c\ne0\)
\item \(\frac ab\cdot\frac cd=\frac{ac}{bd}\)
\item \(\frac ab\div\frac cd=\frac ab\cdot\frac dc\)
\item \(\frac ab+\frac cb=\frac{a+c}{b}\)
\item \(\sqrt{ab}=\sqrt a\sqrt b\)
\item Rationalizing: \(\frac1{\sqrt a}=\frac{\sqrt a}{a}\)
\end{itemize}

\section{Sample Problems}

\subsection{Sample 1}

Simplify \(\frac{x^2+5x}{x}\).

\[\frac{x(x+5)}{x}=x+5,\quad x\ne0\]

\subsection{Sample 2}

Simplify \(\sqrt{72}\).

\[\sqrt{72}=\sqrt{36\cdot2}=6\sqrt2\]

\subsection{Sample 3}

Find the inverse of \(f(x)=\sqrt{x-4}\).

\[y=\sqrt{x-4}\]
\[x=\sqrt{y-4}\]
\[x^2=y-4\]
\[y=x^2+4\]

Answer: \(f^{-1}(x)=x^2+4\), with inverse domain \(x\ge0\).

\section{Classwork}

\begin{enumerate}
\item Simplify \(\frac{6x}{3}\).
\item Simplify \(\frac{x^2+3x}{x}\).
\item Simplify \(\frac{x^2-9}{x-3}\).
\item Multiply \(\frac{2x}{5}\cdot\frac{15}{4x}\).
\item Divide \(\frac{x}{3}\div\frac{2x}{9}\).
\item Add \(\frac{3}{x}+\frac{5}{x}\).
\item Simplify \(\sqrt{50}\).
\item Simplify \(\sqrt{27}\).
\item Rationalize \(\frac{4}{\sqrt2}\).
\item Find the inverse of \(f(x)=x-9\).
\end{enumerate}

\section{Homework}

\begin{enumerate}
\item Simplify \(\frac{10x}{5}\).
\item Simplify \(\frac{x^2-4x}{x}\).
\item Simplify \(\frac{x^2-16}{x-4}\).
\item Multiply \(\frac{3x}{7}\cdot\frac{14}{9x}\).
\item Divide \(\frac{2x}{5}\div\frac{x}{10}\).
\item Add \(\frac{7}{x}+\frac{2}{x}\).
\item Simplify \(\sqrt{98}\).
\item Simplify \(\sqrt{48}\).
\item Rationalize \(\frac{3}{\sqrt3}\).
\item Find the inverse of \(f(x)=x+11\).
\end{enumerate}

\section{Teacher Answer Key}

\subsection{Classwork Answers}

\begin{enumerate}
\item \(2x\)
\item \(x+3\), \(x\ne0\)
\item \(x+3\), \(x\ne3\)
\item \(\frac32\)
\item \(\frac32\)
\item \(\frac8x\)
\item \(5\sqrt2\)
\item \(3\sqrt3\)
\item \(2\sqrt2\)
\item \(f^{-1}(x)=x+9\)
\end{enumerate}

\subsection{Homework Answers}

\begin{enumerate}
\item \(2x\)
\item \(x-4\), \(x\ne0\)
\item \(x+4\), \(x\ne4\)
\item \(\frac23\)
\item \(4\)
\item \(\frac9x\)
\item \(7\sqrt2\)
\item \(4\sqrt3\)
\item \(\sqrt3\)
\item \(f^{-1}(x)=x-11\)
\end{enumerate}

\bigskip\hrule\bigskip

\chapter{Unit 12: Statistics, Data Displays, and Linear Models}

\section{Learning Targets}

\begin{itemize}
\item Calculate measures of center and spread.
\item Interpret dot plots, histograms, and box plot values computationally.
\item Build and use two-way frequency tables.
\item Calculate relative frequencies.
\item Use linear models, residuals, and correlation values.
\end{itemize}

\section{Unit Formula Sheet}

\begin{itemize}
\item Mean: \(\bar{x}=\frac{\sum x}{n}\)
\item Median: middle number in ordered list.
\item Mode: most frequent value.
\item Range: \(\max-\min\)
\item IQR: \(Q_3-Q_1\)
\item Residual: \(y-\hat y\)
\item Linear model prediction: substitute \(x\) into \(\hat y=mx+b\)
\end{itemize}

\section{Sample Problems}

\subsection{Sample 1}

Find the mean of \(4,6,8,10\).

\[\bar{x}=\frac{4+6+8+10}{4}=7\]

\subsection{Sample 2}

Find the median of \(3,9,4,10,6\).

Order: \(3,4,6,9,10\). Median: \(6\)

\subsection{Sample 3}

For \(\hat y=2x+1\), find the residual when \(x=4\) and observed \(y=12\).

Predicted:
\[\hat y=2(4)+1=9\]
Residual:
\[12-9=3\]

\section{Classwork}

\begin{enumerate}
\item Find the mean of \(2,4,6,8\).
\item Find the median of \(7,1,9,3,5\).
\item Find the mode of \(2,3,3,4,5,5,5\).
\item Find the range of \(12,4,18,9\).
\item Find \(Q_1\), \(Q_3\), and IQR for \(2,4,6,8,10,12,14\).
\item For \(\hat y=3x-2\), predict \(y\) when \(x=5\).
\item For \(\hat y=4x+1\), find the residual when \(x=2\) and observed \(y=12\).
\item Compute the joint relative frequency: table total \(40\), cell count \(10\).
\item Compute the marginal relative frequency: table total \(50\), row total \(20\).
\item State whether \(r=0.92\) shows positive or negative correlation.
\end{enumerate}

\section{Homework}

\begin{enumerate}
\item Find the mean of \(5,10,15,20\).
\item Find the median of \(8,2,12,4,10,6\).
\item Find the mode of \(1,1,2,3,3,3,4\).
\item Find the range of \(21,15,7,30\).
\item Find \(Q_1\), \(Q_3\), and IQR for \(1,3,5,7,9,11,13\).
\item For \(\hat y=2x+6\), predict \(y\) when \(x=7\).
\item For \(\hat y=5x-3\), find the residual when \(x=4\) and observed \(y=20\).
\item Compute the joint relative frequency: table total \(80\), cell count \(12\).
\item Compute the marginal relative frequency: table total \(60\), column total \(15\).
\item State whether \(r=-0.81\) shows positive or negative correlation.
\end{enumerate}

\section{Teacher Answer Key}

\subsection{Classwork Answers}

\begin{enumerate}
\item \(5\)
\item \(5\)
\item \(5\)
\item \(14\)
\item \(Q_1=4,\ Q_3=12,\ IQR=8\)
\item \(13\)
\item \(3\)
\item \(0.25\) or \(25\%\)
\item \(0.4\) or \(40\%\)
\item positive
\end{enumerate}

\subsection{Homework Answers}

\begin{enumerate}
\item \(12.5\)
\item \(7\)
\item \(3\)
\item \(23\)
\item \(Q_1=3,\ Q_3=11,\ IQR=8\)
\item \(20\)
\item \(3\)
\item \(0.15\) or \(15\%\)
\item \(0.25\) or \(25\%\)
\item negative
\end{enumerate}

\bigskip\hrule\bigskip

\chapter{Unit 13: Cumulative Review and Final Exam Practice}

\section{Cumulative Formula Sheet}

\begin{itemize}
\item Exponents: \(a^ma^n=a^{m+n}\), \((a^m)^n=a^{mn}\), \(a^{-n}=\frac1{a^n}\)
\item Radicals: \(a^{1/n}=\sqrt[n]a\)
\item Slope: \(m=\frac{y_2-y_1}{x_2-x_1}\)
\item Line: \(y=mx+b\)
\item Quadratic formula: \(x=\frac{-b\pm\sqrt{b^2-4ac}}{2a}\)
\item Axis: \(x=-\frac b{2a}\)
\item Average rate of change: \(\frac{f(b)-f(a)}{b-a}\)
\item Arithmetic sequence: \(a_n=a_1+(n-1)d\)
\item Geometric sequence: \(a_n=a_1r^{n-1}\)
\item Mean: \(\bar{x}=\frac{\sum x}{n}\)
\end{itemize}

\section{Final Exam Practice A: Student Version}

\begin{enumerate}
\item Simplify \(3(2x-5)+4x\).
\item Simplify \(\frac{x^9}{x^4}\).
\item Simplify \(64^{2/3}\).
\item Solve \(5x-7=18\).
\item Solve \(-2x+3\le11\).
\item Solve \(|x-6|=4\).
\item Find the slope through \((2,3)\) and \((8,15)\).
\item Write the equation of a line with slope \(-3\) and \(y\)-intercept \(5\).
\item Solve the system \(y=x+2\), \(y=-x+8\).
\item Evaluate \(f(x)=x^2-2x\) at \(x=5\).
\item Find the average rate of change of \(f(x)=3x+1\) from \(x=0\) to \(x=4\).
\item Multiply \((x+7)(x-2)\).
\item Factor \(x^2+2x-15\).
\item Solve \(x^2-4x-12=0\).
\item Find the vertex of \(y=x^2-10x+9\).
\item Determine growth or decay: \(y=4(0.75)^x\).
\item Write an explicit formula for \(3,8,13,18,\ldots\).
\item Write an explicit formula for \(2,6,18,54,\ldots\).
\item Simplify \(\sqrt{75}\).
\item Find the mean of \(6,8,10,12\).
\end{enumerate}

\section{Final Exam Practice B: Student Version}

\begin{enumerate}
\item Simplify \(-4(x+3)+9x\).
\item Simplify \((a^3)^4\).
\item Simplify \(125^{1/3}\).
\item Solve \(\frac{x}{4}+6=10\).
\item Solve \(3x-8>13\).
\item Solve \(|2x+1|=7\).
\item Find the slope through \((-1,4)\) and \((3,12)\).
\item Write the equation of a line with slope \(\frac23\) and \(y\)-intercept \(-1\).
\item Solve the system \(x+y=11\), \(x-y=3\).
\item Evaluate \(g(x)=2x^2+1\) at \(x=-3\).
\item Find the average rate of change of \(f(x)=x^2\) from \(x=2\) to \(x=6\).
\item Multiply \((x-5)(x-4)\).
\item Factor \(x^2-11x+30\).
\item Solve \(x^2+9x+20=0\).
\item Find the vertex of \(y=x^2+6x-1\).
\item Determine growth or decay: \(y=2(1.4)^x\).
\item Write an explicit formula for \(10,7,4,1,\ldots\).
\item Write an explicit formula for \(5,20,80,320,\ldots\).
\item Simplify \(\sqrt{108}\).
\item Find the median of \(9,1,5,7,3\).
\end{enumerate}

\section{Teacher Answer Key}

\subsection{Final Exam Practice A Answers}

\begin{enumerate}
\item \(10x-15\)
\item \(x^5\)
\item \(16\)
\item \(x=5\)
\item \(x\ge-4\)
\item \(x=10\) or \(x=2\)
\item \(2\)
\item \(y=-3x+5\)
\item \((3,5)\)
\item \(15\)
\item \(3\)
\item \(x^2+5x-14\)
\item \((x+5)(x-3)\)
\item \(x=6,-2\)
\item \((5,-16)\)
\item decay
\item \(a_n=3+5(n-1)\)
\item \(a_n=2(3)^{n-1}\)
\item \(5\sqrt3\)
\item \(9\)
\end{enumerate}

\subsection{Final Exam Practice B Answers}

\begin{enumerate}
\item \(5x-12\)
\item \(a^{12}\)
\item \(5\)
\item \(x=16\)
\item \(x>7\)
\item \(x=3\) or \(x=-4\)
\item \(2\)
\item \(y=\frac23x-1\)
\item \((7,4)\)
\item \(19\)
\item \(8\)
\item \(x^2-9x+20\)
\item \((x-5)(x-6)\)
\item \(x=-4,-5\)
\item \((-3,-10)\)
\item growth
\item \(a_n=10-3(n-1)\)
\item \(a_n=5(4)^{n-1}\)
\item \(6\sqrt3\)
\item \(5\)
\end{enumerate}

\bigskip\hrule\bigskip

\chapter{Teacher Implementation Notes}

\section{Suggested Weekly Routine}

\begin{itemize}
\item \textbf{Day 1:} Formula sheet, vocabulary, and sample problems.
\item \textbf{Day 2:} Guided classwork set, error analysis, and fluency drills.
\item \textbf{Day 3:} Mixed practice and short quiz.
\item \textbf{Day 4:} Homework review, reteaching, and extension problems.
\item \textbf{Day 5:} Unit checkpoint or cumulative review.
\end{itemize}

\section{Computational Practice Expectations}

Students should show:

\begin{itemize}
\item Substitution steps for evaluation problems.
\item Inverse operations for equations.
\item Sign changes when solving inequalities involving multiplication or division by negatives.
\item Factoring steps for quadratic equations.
\item Formula substitution for slope, vertex, quadratic formula, statistics, and sequences.
\end{itemize}

\section{Recommended Assessment Types}

\begin{itemize}
\item 5-question daily entrance ticket.
\item 10-question weekly computational quiz.
\item Unit test with 20 to 30 computational items.
\item Cumulative final exam using Units 1 through 13.
\end{itemize}

\end{document}
